\documentclass[12pt]{article}
\usepackage{amsmath, amssymb, amsfonts, geometry, setspace, tcolorbox, amsthm}
\geometry{margin=1in}
\setstretch{1.3}
\setlength{\parindent}{0pt}

\tcbset{colback=gray!5, colframe=black!70, boxrule=0.6pt, arc=2pt, 
left=6pt, right=6pt, top=4pt, bottom=4pt}

\newtheorem{definition}{Definition}
\newtheorem{theorem}{Theorem}
\newtheorem{example}{Example}

\title{\textbf{Integration and Economic Dynamics}\\[7pt]
\large ECON 320 - Introduction to Mathematical Economics}
\author{Muhebullah Karimzada}
\date{Fall 2025}

\begin{document}
\maketitle





%=========================================================
\section*{1. Motivation: Why Integration in Economics?}

So far, you may be familiar with \textbf{derivatives}: 
\[
\frac{dy}{dx}, \quad f'(x), \quad \frac{dQ}{dP}, \quad \text{etc.}
\]
Derivatives describe \emph{instantaneous change}. For example:
\begin{itemize}
\item Marginal cost is the derivative of total cost with respect to quantity.
\item Marginal utility is the derivative of total utility with respect to consumption.
\item Marginal product is the derivative of output with respect to an input.
\end{itemize}

However, many economic questions are about \textbf{accumulation over time or over quantity}, such as:
\begin{itemize}
\item Given marginal cost, what is total cost between $Q=0$ and $Q=100$?
\item Given investment as a function of time, what will the capital stock be after 10 years?
\item Given an income flow $Y(t)$, what is the present value of all future income?
\end{itemize}

To answer such questions, we need a mathematical tool that \emph{adds up} small contributions:
\[
\text{Total} = \text{sum of many tiny pieces}.
\]
This is exactly what \textbf{integration} does.



%=========================================================
\section*{2. What is an Integral? First Intuition}

\subsection*{2.1 Area Under a Curve as a Sum}

Consider a function $f(x)$ defined on the interval $[a,b]$. Visually, we can think of $f(x)$ as a curve. Suppose we want the \textbf{area under the curve} from $x=a$ to $x=b$.

A natural idea:
\begin{enumerate}
\item Divide $[a,b]$ into many small sub-intervals of width $\Delta x$.
\item On each small interval, approximate the area by a rectangle:
\[
\text{height} \approx f(x_i), \quad \text{width} = \Delta x.
\]
\item Add up all these small rectangular areas:
\[
\text{Total area} \approx \sum_i f(x_i)\Delta x.
\]
\end{enumerate}
As we make $\Delta x$ smaller and smaller (more and more rectangles), this approximation becomes better and better.

\subsection*{2.2 The Limit: From Sum to Integral}

Formally, the \textbf{definite integral} of $f(x)$ from $a$ to $b$ is defined as:
\[
\int_a^b f(x)\,dx = \lim_{\Delta x \to 0} \sum_i f(x_i)\Delta x.
\]

This is a limit of sums. That is why we say the integral is a \emph{continuous sum}.

\begin{tcolorbox}
\textbf{Definite Integral:} $\displaystyle \int_a^b f(x)\,dx$ is the limiting value of sums of the form $\sum f(x_i)\Delta x$ as $\Delta x\to 0$. It represents the \emph{area under the curve} $y=f(x)$ from $x=a$ to $x=b$ (when $f(x)\ge 0$).
\end{tcolorbox}
\newpage
%=========================================================
\section*{3. Relationship Between Derivatives and Integrals}

It turns out that there is a deep relationship between differentiation and integration.

Suppose $F(x)$ is a function such that 
\[
F'(x) = f(x).
\]
Then we say $F(x)$ is an \textbf{antiderivative} (or \textbf{indefinite integral}) of $f(x)$.

\begin{tcolorbox}
\textbf{Fundamental Theorem of Calculus (informal)}  

If $F'(x) = f(x)$, then:
\[
\int_a^b f(x)\, dx = F(b) - F(a).
\]

In words: to compute the definite integral of $f(x)$ from $a$ to $b$, find any function $F(x)$ whose derivative is $f(x)$, then take $F(b)-F(a)$.
\end{tcolorbox}

This theorem tells us:
\begin{itemize}
\item Integration can be done by finding antiderivatives.
\item Differentiation and integration are inverse operations.
\end{itemize}

%=========================================================
\section*{4. Indefinite Integrals: Basic Rules (With Steps)}

An \textbf{indefinite integral} has the form:
\[
\int f(x)\, dx = F(x) + C,
\]
where $F'(x) = f(x)$ and $C$ is an arbitrary constant (because the derivative of a constant is zero).

\subsection*{4.1 Power Rule}

We start from the derivative:
\[
\frac{d}{dx}x^{n+1} = (n+1)x^n.
\]
We want to \emph{reverse} this. Suppose we want $\int x^n dx$.

\textbf{Step 1:} Guess that the antiderivative has the form $A x^{n+1}$.

\textbf{Step 2:} Differentiate:
\[
\frac{d}{dx} (A x^{n+1}) = A (n+1) x^n.
\]

\textbf{Step 3:} We want this to equal $x^n$, so set:
\[
A (n+1) = 1 \quad \Rightarrow \quad A = \frac{1}{n+1}.
\]

\textbf{Conclusion:}
\[
\int x^n dx = \frac{x^{n+1}}{n+1} + C, \quad \text{for } n\ne -1.
\]

\subsection*{4.2 Integral of $1/x$}

We know 
\[
\frac{d}{dx} \ln|x| = \frac{1}{x}.
\]

Thus:
\[
\int \frac{1}{x} dx = \ln|x| + C.
\]

\subsection*{4.3 Exponential Functions}

We know:
\[
\frac{d}{dx} e^{ax} = a e^{ax}.
\]

To find $\int e^{ax}dx$, follow similar steps:

\textbf{Step 1:} Guess that the antiderivative is $B e^{ax}$.

\textbf{Step 2:} Differentiate:
\[
\frac{d}{dx} (B e^{ax}) = B a e^{ax}.
\]

\textbf{Step 3:} We want this to equal $e^{ax}$, so:
\[
Ba = 1 \quad \Rightarrow \quad B = \frac{1}{a}.
\]

\textbf{Conclusion:}
\[
\int e^{ax} dx = \frac{1}{a} e^{ax} + C.
\]

\subsection*{4.4 Linearity of the Integral}

If $k$ is a constant and $f,g$ are functions:
\[
\int [kf(x)]dx = k\int f(x)dx,
\]
\[
\int [f(x) + g(x)]dx = \int f(x)dx + \int g(x)dx.
\]

We can justify this by using the linearity of differentiation and then reversing the process.

\textbf{Example (step-by-step):}
\[
\int (3x^2 + 4x)dx.
\]

\textbf{Step 1:} Apply linearity:
\[
\int (3x^2 + 4x)dx = \int 3x^2 dx + \int 4x dx.
\]

\textbf{Step 2:} Pull constants out:
\[
=3 \int x^2 dx + 4\int x dx.
\]

\textbf{Step 3:} Use power rule:
\[
=3 \left(\frac{x^3}{3}\right) + 4\left(\frac{x^2}{2}\right) + C
= x^3 + 2x^2 + C.
\]
\newpage
%=========================================================
\section*{5. Definite Integrals and the Fundamental Theorem of Calculus}

\subsection*{5.1 Definite Integrals}

A \textbf{definite integral} from $a$ to $b$ is written:
\[
\int_a^b f(x)dx.
\]

This is a number (not a function). Using an antiderivative $F(x)$ with $F'(x)=f(x)$:

\[
\int_a^b f(x)dx = F(b) - F(a).
\]

\textbf{Example:}
\[
\int_0^2 (3x^2 + 4x)dx.
\]

We already know antiderivative:
\[
\int (3x^2 + 4x)dx = x^3 + 2x^2 + C.
\]

So:
\[
\int_0^2 (3x^2 + 4x)dx = [x^3 + 2x^2]_0^2 = (2^3 + 2\cdot 2^2) - (0^3 + 2\cdot 0^2).
\]

Compute step-by-step:
\[
2^3 = 8, \quad 2^2 = 4, \quad 2\cdot 4 = 8.
\]
So:
\[
x^3 + 2x^2 \text{ at } x=2 \text{ is } 8 + 8 = 16,
\]
and at $x=0$:
\[
0^3 + 2\cdot 0^2 = 0.
\]

Therefore:
\[
\int_0^2 (3x^2 + 4x)dx = 16 - 0 = 16.
\]

\begin{tcolorbox}
\textbf{Economic Reading:} If $3x^2+4x$ were the marginal cost, then the total cost of producing from 0 to 2 units would be 16 (plus any fixed cost).
\end{tcolorbox}

%=========================================================
\section*{6. Economic Applications: Static (No Time Yet)}

\subsection*{6.1 Total Cost from Marginal Cost}

Suppose marginal cost is $MC(Q)$, and fixed cost is $F$. Then:
\[
TC(Q) = F + \int_0^Q MC(q)dq.
\]

\textbf{Example:}
\[
MC(Q) = 10 + 2Q, \quad F=50.
\]

\textbf{Step 1:} Compute the integral:
\[
\int_0^Q (10+2q)dq = \int_0^Q 10\,dq + \int_0^Q 2q\,dq.
\]

\textbf{Step 2:} Integrate separately:
\[
\int_0^Q 10\,dq = 10Q, \quad
\int_0^Q 2q\,dq = 2\cdot \frac{Q^2}{2} = Q^2.
\]

\textbf{Step 3:} Add:
\[
\int_0^Q (10+2q)dq = 10Q + Q^2.
\]

\textbf{Step 4:} Add fixed cost:
\[
TC(Q) = 50 + 10Q + Q^2.
\]

\subsection*{6.2 Consumer Surplus}

Demand curve: $P(Q)$, market price: $P^*$. Quantity traded: $Q^*$.

\[
CS = \int_0^{Q^*} P(Q)dQ - P^* Q^*.
\]

This is literally the area under the demand curve above the horizontal price line.
\newpage
%=========================================================
\section*{7. Integration and Economic Dynamics: Stocks and Flows}

Now we move from \emph{static} applications to \emph{dynamic} ones.

\subsection*{7.1 Flows Accumulating into Stocks}

Let $I(t)$ be net investment (a \textbf{flow}), measured per unit of time.  

Capital stock $K(t)$ is a \textbf{stock}.

The change in capital over a small time $\Delta t$ is approximately:
\[
\Delta K \approx I(t)\Delta t.
\]

To get total change over $[0,t]$, we add up all tiny changes:
\[
K(t) - K(0) \approx \sum I(t_i)\Delta t.
\]

In the limit:
\[
K(t) - K(0) = \int_0^t I(s)ds.
\]

So:
\[
\boxed{K(t) = K(0) + \int_0^t I(s)ds.}
\]

This is the basic \textbf{stock-flow} relationship.

\subsection*{7.2 Example: Constant Investment}

Suppose $I(t) = \bar{I}$, a constant.

Then:
\[
K(t) = K(0) + \int_0^t \bar{I}\ ds = K(0) + \bar{I} t.
\]

This is just a straight line: capital grows at a constant rate over time.
\newpage
%=========================================================
\section*{8. Introducing Differential Equations}

A \textbf{differential equation} is an equation that involves an unknown function and its derivatives.

Example:
\[
\frac{dK}{dt} = I(t).
\]

This equation says: the rate of change of $K$ with respect to time equals the investment flow.

If we know $I(t)$, we can integrate to get $K(t)$:
\[
K(t)=K(0)+\int_0^t I(s)ds.
\]

So integration is the natural tool for solving many differential equations.

%=========================================================
\section*{9. Example: Capital Accumulation with Depreciation}

Consider:
\[
\frac{dK}{dt} = I - \delta K,
\]
where:
\begin{itemize}
\item $I$ is \emph{constant net investment}.
\item $\delta$ is the \emph{depreciation rate}.
\end{itemize}

This is a first-order linear differential equation.

\subsection*{9.1 Rewrite the Equation}

We can rewrite as:
\[
\frac{dK}{dt} + \delta K = I.
\]

Now it is in the standard form:
\[
\frac{dK}{dt} + P(t)K = Q(t),
\]
with $P(t)=\delta$ (constant), $Q(t)=I$ (constant).

\subsection*{9.2 Integrating Factor Method (Step by Step)}

We choose an \textbf{integrating factor} $\mu(t)$ such that:
\[
\mu(t) = e^{\int P(t)dt} = e^{\int \delta dt} = e^{\delta t}.
\]

\textbf{Step 1:} Multiply both sides of the differential equation by $\mu(t) = e^{\delta t}$:
\[
e^{\delta t}\frac{dK}{dt} + \delta e^{\delta t}K = I e^{\delta t}.
\]

\textbf{Step 2:} Recognize that the left-hand side is:
\[
\frac{d}{dt}(e^{\delta t}K).
\]
We can verify this:

Take $e^{\delta t}K$ and differentiate with respect to $t$:
\[
\frac{d}{dt}(e^{\delta t}K) = e^{\delta t}\frac{dK}{dt} + \delta e^{\delta t}K.
\]
So indeed:
\[
\frac{d}{dt}(e^{\delta t}K) = I e^{\delta t}.
\]

\textbf{Step 3:} Integrate both sides from $0$ to $t$:
\[
\int_0^t \frac{d}{ds}(e^{\delta s}K(s)) ds = \int_0^t I e^{\delta s} ds.
\]

Left-hand side:
\[
\int_0^t \frac{d}{ds}(e^{\delta s}K(s)) ds = e^{\delta t}K(t) - e^{0}K(0) = e^{\delta t}K(t) - K(0).
\]

Right-hand side:
\[
\int_0^t I e^{\delta s} ds = I \int_0^t e^{\delta s}ds.
\]
Compute:
\[
\int_0^t e^{\delta s}ds = \left[\frac{1}{\delta} e^{\delta s}\right]_0^t = \frac{1}{\delta} e^{\delta t} - \frac{1}{\delta} e^{0} = \frac{1}{\delta}(e^{\delta t} - 1).
\]

Thus:
\[
\int_0^t I e^{\delta s} ds = I \cdot \frac{1}{\delta}(e^{\delta t} - 1) = \frac{I}{\delta}(e^{\delta t} - 1).
\]

\textbf{Step 4:} Put both sides together:
\[
e^{\delta t}K(t) - K(0) = \frac{I}{\delta}(e^{\delta t} - 1).
\]

\textbf{Step 5:} Solve for $K(t)$:
\[
e^{\delta t}K(t) = K(0) + \frac{I}{\delta}(e^{\delta t} - 1).
\]

Divide by $e^{\delta t}$:
\[
K(t) = K(0)e^{-\delta t} + \frac{I}{\delta}(1 - e^{-\delta t}).
\]

\begin{tcolorbox}
\textbf{Capital Accumulation Solution}  
\[
K(t) = K(0)e^{-\delta t} + \frac{I}{\delta}(1 - e^{-\delta t}).
\]

\begin{itemize}
\item $K(0)e^{-\delta t}$ is the depreciating initial stock.
\item $\frac{I}{\delta}(1-e^{-\delta t})$ captures the build-up of new capital from investment.
\item As $t\to\infty$, $e^{-\delta t}\to 0$, so
\[
K(\infty) = \frac{I}{\delta},
\]
which is the steady-state capital stock where net investment equals depreciation.
\end{itemize}
\end{tcolorbox}

\newpage
%=========================================================
\section*{10. Debt Dynamics: Another Example}

Let $B(t)$ be government debt at time $t$. Suppose:
\[
\frac{dB}{dt} = rB + D(t),
\]
where:
\begin{itemize}
\item $r$ is the interest rate on debt.
\item $D(t)$ is the (primary) deficit at time $t$.
\end{itemize}

\subsection*{10.1 Rearranging and Integrating Factor}

Rewrite:
\[
\frac{dB}{dt} - rB = D(t).
\]

Here $P(t) = -r$, so:
\[
\mu(t) = e^{\int -r dt} = e^{-rt}.
\]

Multiply both sides by $e^{-rt}$:
\[
e^{-rt}\frac{dB}{dt} - r e^{-rt} B = D(t) e^{-rt}.
\]

Left side is:
\[
\frac{d}{dt}(e^{-rt}B).
\]

Check:
\[
\frac{d}{dt}(e^{-rt}B) = e^{-rt}\frac{dB}{dt} + (-r e^{-rt})B = e^{-rt}\frac{dB}{dt} - r e^{-rt} B.
\]
Correct.

So we have:
\[
\frac{d}{dt}(e^{-rt}B) = D(t) e^{-rt}.
\]

Integrate from $0$ to $t$:
\[
\int_0^t \frac{d}{ds}(e^{-rs}B(s)) ds = \int_0^t D(s)e^{-rs}ds.
\]

Left-hand side:
\[
e^{-rt}B(t) - e^{0}B(0) = e^{-rt}B(t) - B(0).
\]

Thus:
\[
e^{-rt}B(t) - B(0) = \int_0^t D(s)e^{-rs}ds.
\]

Solve for $B(t)$:
\[
e^{-rt}B(t) = B(0) + \int_0^t D(s)e^{-rs}ds,
\]
\[
B(t) = e^{rt}\left[ B(0) + \int_0^t D(s)e^{-rs}ds \right].
\]

\begin{tcolorbox}
\textbf{Debt Solution}  
\[
B(t) = e^{rt}\left[ B(0) + \int_0^t D(s)e^{-rs}ds \right].
\]

Interpretation:
\begin{itemize}
\item $B(0)e^{rt}$ is the initial debt grown at interest rate $r$.
\item The integral term represents the accumulation of discounted deficits over time.
\end{itemize}
\end{tcolorbox}

%=========================================================
\section*{11. Present Value as an Integral}

In continuous time, if $Y(t)$ is a flow of income and $r$ is the constant discount rate, the present value (PV) is:
\[
PV = \int_0^\infty Y(t) e^{-rt}dt.
\]

\subsection*{11.1 Example: Constant Income Stream}

Suppose $Y(t)=Y$ (constant). Then:
\[
PV = \int_0^\infty Y e^{-rt}dt = Y\int_0^\infty e^{-rt}dt.
\]

Compute the integral:
\[
\int_0^\infty e^{-rt}dt = \lim_{T\to\infty} \int_0^T e^{-rt}dt.
\]

First compute up to $T$:
\[
\int_0^T e^{-rt}dt = \left[-\frac{1}{r}e^{-rt}\right]_0^T = -\frac{1}{r}e^{-rT} + \frac{1}{r}e^{0} = \frac{1}{r}(1 - e^{-rT}).
\]

Now take the limit as $T\to\infty$:
\[
\lim_{T\to\infty} \frac{1}{r}(1 - e^{-rT}) = \frac{1}{r}(1-0) = \frac{1}{r}.
\]

So:
\[
PV=Y\cdot \frac{1}{r} = \frac{Y}{r}.
\]

This is the continuous-time analog of the perpetuity formula.

%=========================================================
\section*{12. Summary and Economic Insights}

\begin{tcolorbox}
\textbf{Core Messages}
\begin{itemize}
\item Integration is the mathematical tool that sums infinitely many tiny contributions (areas, flows).
\item The definite integral is the limit of Riemann sums and gives areas under curves.
\item The indefinite integral (antiderivative) reverses differentiation.
\item In economics, integration is used to compute total cost, total revenue, consumer surplus, and more.
\item A stock (capital, debt, knowledge) is often the integral of its net flow (investment, deficit, R\&D).
\item Differential equations model how economic variables evolve over continuous time.
\item The integrating factor method allows us to solve a wide class of linear first-order differential equations.
\item Present value in continuous time is an integral of discounted flows.
\end{itemize}
\end{tcolorbox}


\end{document}
